\section{The Futures Lifecycle}
\label{sec:futures}

\providecommand{\code}[1]{\texttt{#1}}
\providecommand{\ac}{\texttt{ac}}
\providecommand{\netq}{\mathrm{net\_qty}}
\providecommand{\unit}{u_{F}}
\providecommand{\mult}{m}

A listed future stresses every primitive at once. An equity trade is one transaction: cash one
way, shares the other. A futures trade reads the same, ``buy 10 at 100,'' yet moves no cash at
inception. Its economic effect unfolds over days, through variation margin, partial closes, and
settlement. The future is therefore the canonical lifecycle. It is recorded, traded, settled daily,
expired, and flattened to zero. Each step is one atomic \code{FutStateDelta} over the move and
conservation primitives of \S\ref{sec:ledger} and the three-home state model of \S\ref{sec:states}.
The \code{Fut*} types form a futures-engine kernel that \emph{projects onto} the core
\code{Transaction}, exactly as the CDM and settlement boundary types do, under the same
naming discipline. The kernel is not a second core. Conservation is shown at each step, not
asserted. The centrepiece is variation-margin settlement: one event writes the shared settlement
price once, fans out over the current holders, resets each accumulated cost, and moves cash.
That settlement price is an externally sourced mark, consumed under single-basis consumption
(Invariant~\ref{inv:basis}): the mark and the positions it settles stand in one basis.

%==============================================================================
\subsection{Notation, and where the future's state lives}
%==============================================================================

The primitives are those of \S\ref{sec:ledger}: wallet, unit, move, conservation. The homes are
those of \S\ref{sec:states}. Write $\netq(w,u):=h(w,u)$ for the signed holding of wallet $w$
in unit $u$.

\emph{UnitStatus} is a projection, not a source of truth: a shared per-unit read cache of
the event log, rebuilt exactly by replay (\S\ref{sec:states-unitstatus}).

\noindent For the future, the fields divide by who reads them and how often they change:

\begin{center}\small
\begin{tabular}{@{}p{6.2cm}p{3.0cm}p{5.6cm}@{}}
\toprule
\textbf{Field} & \textbf{Home} & \textbf{Reason} \\
\midrule
\code{multiplier}, \code{currency}, \code{expiry}, \code{clearinghouse}, \code{exchange},
\code{product\_id} & \emph{ProductTerms}$[\unit]$ & immutable; the CME and ICE contracts are
distinct units. \\
\addlinespace
\code{lifecycle\_stage}, \code{last\_settlement\_price}, \code{last\_settlement\_date} &
\emph{UnitStatus}$[\unit]$ & one settlement value per contract, read identically by every holder. \\
\addlinespace
\code{accumulated\_cost} (\ac) & \emph{PositionState}$[w,\unit]$ & two wallets holding the same
contract carry different \ac. \\
\bottomrule
\end{tabular}
\end{center}

\noindent Accumulated cost \ac{} is a \emph{conserved} field: $\sum_w \ac(w,u)=0$ for every $u$.
The identity is enforced one level up, at the event handler (C2), because no production language
provides a refinement type over a sum of exact integers for free.

\paragraph{The atomic delta.} One event proposes one \code{FutStateDelta} for one unit across the
three maps: at most one shared \code{FutStage} write, a map of per-holder position deltas
$(\Delta\netq,\Delta\ac)$, and a map of per-holder cash legs. The delta applies in full or not at
all (C3). A partially applied delta is not representable. A delta reaches application only through
\code{futValidate}. That gate discharges conservation --- $\sum_w\Delta f(w,u)=0$ for each conserved
field --- and returns the abstract witness \code{FutValidDelta}. An unconserved delta cannot be applied.

\noindent\emph{The futures kernel is a boundary, not a core.} \code{FutStateDelta} and its validated
form \code{FutValidDelta} are the futures engine's own constructions, named apart from the core
\code{Transaction}. Each validated futures delta \emph{projects onto} one core
\code{Transaction}, recorded through the single core door \code{applyTx}. The projection carries
the conserving moves (the edge-sum of the per-holder legs) together with the non-balance rows (the
\ac{} resets and the shared stage write). The local \code{futValidate} discharges conservation
\emph{within} the futures kernel before it projects out; the core therefore keeps no separate
validate gate and remains sealed.

\paragraph{The single dimension bridge.} Contracts, money, and price are distinct types
(\code{Qty}, \code{Cash}, \code{Price}). The only multiplication in the lifecycle,
\[
\mathrm{futMark}(\netq, S, \mult) \;=\; \netq \cdot S \cdot \mult \;:\; \code{Cash},
\]
converts contracts at a price into money. \code{Price} carries no addition: two prices are never
summed, nor moved between wallets. The load-bearing identity $\mathrm{VM}=\netq\cdot S\cdot\mult+\ac$
typechecks only because both summands are \code{Cash}. \code{Price} and \code{Cash} share one fixed
minor-unit scale, chosen so $\netq\cdot S\cdot\mult$ lands in \code{Cash}'s unit. The worked example
carries \code{Price} in whole index points and \code{Cash} in whole dollars for readability. A
production scale carries \code{Price} in scaled minor units --- a quarter-point for ES, a cent for a
bond quoted in 32nds --- so sub-point ticks are exact. The arithmetic is integer throughout and never
divides, so the choice of scale is free: at any scale there is no rounding to break conservation.

The scalar layer --- the three dimensions, \code{futMark}, and the \code{PosQty} parse boundary of
\S\ref{sec:trade} --- is the kernel's foundation. The type discipline encodes one mental model:
quantities and cash add, prices do not, so ``a price summed into a balance'' is a type error rather
than a bug to catch.

\begin{lstlisting}[language=Haskell]
-- Three dimensions kept apart. Qty and Cash are additive groups; Price is not
-- (you never sum two prices, nor move one between wallets) -- so Price has no
-- Monoid instance, and that absence makes "a price summed into a balance"
-- unspellable rather than merely unintended.
newtype Qty   = Qty   Integer deriving (Eq, Ord, Show)   -- additive group: the (<>)/mempty
newtype Cash  = Cash  Integer deriving (Eq, Ord, Show)   --   instances of sec:contracts
newtype Price = Price Integer deriving (Eq, Ord, Show)   -- deliberately NO Monoid
newtype Day   = Day   Int     deriving (Eq, Ord, Show)

qneg    :: Qty  -> Qty ;  qneg    (Qty  n) = Qty  (negate n)
cashNeg :: Cash -> Cash;  cashNeg (Cash n) = Cash (negate n)

-- The single dimension bridge: contracts at a price, scaled by the multiplier,
-- become money. The only multiplication in the lifecycle, and the only place the
-- dimensions cross. VM = futMark ... <> ac typechecks only because both are Cash.
futMark :: Qty -> Price -> Integer -> Cash
futMark (Qty q) (Price s) mult = Cash (q * s * mult)

newtype WalletId = WalletId String deriving (Eq, Ord, Show)
newtype UnitId   = UnitId   String deriving (Eq, Ord, Show)

-- The positive-quantity parse boundary (G5). PosQty is abstract: mkPosQty is its
-- sole constructor, so q <= 0 is unrepresentable downstream of the boundary.
newtype PosQty = PosQty Qty deriving (Eq, Ord, Show)

mkPosQty :: Qty -> Maybe PosQty
mkPosQty q@(Qty n) | n > 0     = Just (PosQty q)
                   | otherwise = Nothing

unPosQty :: PosQty -> Qty
unPosQty (PosQty q) = q
\end{lstlisting}

\noindent\emph{Evaluations.} \code{futMark (Qty 10) (Price 100) 50 = Cash 50000};
\code{futMark (Qty (-10)) (Price 102) 50 <> Cash 51000 = Cash 0} (a \code{Cash}
sum); \code{Price 100 <> Price 102} is a type error --- \code{Price} has no
\code{<>}.

%==============================================================================
\subsection{The instrument and the worked example}
%==============================================================================

\begin{definition}[The worked instrument]
One cash-settled listed future $\unit=\code{ES-FUT}$, with $\mult=50$, currency USD, expiry Day 3,
clearinghouse CME-CH, exchange CME. Wallets A, B, C are traders; CH is the clearinghouse virtual
wallet, the variation-margin settlement hub whose cash leg is the residual that balances the holders'
legs.
\end{definition}

CH stands in as a settlement hub, not as a \emph{novating} central counterparty. Faithful novation
would interpose CH on the \emph{contract} leg of every trade ($A\,{+}q\,/\,\text{CH}\,{-}q$ and
$\text{CH}\,{+}q\,/\,B\,{-}q$), so that the two cleared sides never face each other. Clearing
infrastructure is out of scope. Here CH carries only the variation-margin cash residual. Two rules
govern every figure:

\begin{itemize}
\item \textbf{Trade.} For a signed-quantity change $\Delta$ at price $p$, each leg sets
$\ac \mathrel{+}= -\,\Delta\cdot p\cdot \mult$.
\item \textbf{Settlement.} At price $S$, the per-holder variation margin is
$\mathrm{VM}(w)=\netq(w)\cdot S\cdot \mult + \ac(w)$ (cash, routed through CH); then
$\ac(w)\leftarrow -\,\netq(w)\cdot S\cdot \mult$.
\end{itemize}

\noindent The whole life, with $\netq$, \ac, and per-step variation-margin cash for $(A,B,C)$:

\begin{center}
\resizebox{\textwidth}{!}{%
\begin{tabular}{@{}llrrr@{}}
\toprule
\textbf{Step} & \textbf{Event} & \textbf{$\netq$ (A,B,C)} & \textbf{\ac{} (A,B,C)} & \textbf{VM cash (A,B,C)} \\
\midrule
Listing    & $\unit$ registered; stage \code{REGISTERED}, no mark; no rows & $(0,0,0)$ & $(-,-,-)$ & $-$ \\
T1         & A buys 10 from B @100                       & $(10,-10,0)$ & $(-50000,+50000,0)$ & $-$ \\
Settle d1  & $S_1=102$; shared price $\leftarrow 102$    & $(10,-10,0)$ & $(-51000,+51000,0)$ & $(+1000,-1000,0)$ \\
T2         & C buys 4 from A @103                        & $(6,-10,4)$  & $(-30400,+51000,-20600)$ & $-$ \\
Settle d2  & $S_2=101$; shared price $\leftarrow 101$    & $(6,-10,4)$  & $(-30300,+50500,-20200)$ & $(-100,+500,-400)$ \\
T3         & B buys 4 from C @101 (C $\to$ flat)         & $(6,-6,0)$   & $(-30300,+30300,0)$ & $-$ \\
Expiry     & stage $\to$ \code{EXPIRED}; final $S=105$   & $(6,-6,0)$   & $(-31500,+31500,0)$ & $(+1200,-1200,0)$ \\
Close      & units returned to CH; positions $\to 0$     & $(0,0,0)$    & $(0,0,0)$ & $-$ \\
\bottomrule
\end{tabular}}
\end{center}

\noindent The subsections below walk this table in event order. Each event is one \code{FutStateDelta}
showing the three conservation sums $\sum\Delta\netq$, $\sum\Delta\ac$, $\sum\mathrm{VM}$. The example
is the argument; the prose names the mechanism.

%==============================================================================
\subsection{Listing}
%==============================================================================

Registration writes \emph{ProductTerms}$[\unit]$ (the six immutable terms) and
\emph{UnitStatus}$[\unit]$ together, establishing ``$u\in$ terms $\iff u\in$ status'' by construction.
The status initialises to stage \code{REGISTERED} with no settlement mark. No \emph{PositionState} row
exists: $\code{futPosition}(w,\unit)=\mathrm{None}$ for every $w$.

\code{REGISTERED} denotes a unit recorded in and known to the ledger before any settlement, not
admission to an exchange. An OTC option or an internal strategy reaches it identically. Two status
states are unreachable: a never-traded unit carrying a price, and an expired unit lacking a final
mark. Both are made unspellable by fusing the mark onto the stage:
\[
\code{FutStage}=\code{FutRegistered}\;\mid\;\code{FutActive}(\code{Maybe Settlement})\;\mid\;
\code{FutExpired}\;\code{Settlement}.
\]
\code{FutRegistered} carries no mark; \code{FutExpired} always carries one. \code{last\_\allowbreak settlement\_\allowbreak price}
and \code{last\_\allowbreak settlement\_\allowbreak date} are projections of the \code{Settlement} carried by the stage, not
independent fields. Conservation holds vacuously: no holders, so every sum is the empty sum, zero
(C9). Re-registering an existing unit id is a hard error, never a silent reset (C10).

The shared status type and its registration follow.

\begin{lstlisting}[language=Haskell]
-- FutTerms: the six immutable contract terms, one per unit, co-registered with
-- FutStatus. ptMultiplier drives mark value and the settlement fan-out. The Fut*
-- prefix marks this engine as a futures-boundary construction, distinct from the
-- core types: it PROJECTS onto the core Transaction (see "The atomic delta").
data FutTerms = FutTerms
  { ptMultiplier    :: !Integer
  , ptCurrency      :: !String
  , ptExpiry        :: !Day
  , ptClearinghouse :: !String
  , ptExchange      :: !String
  , ptProductId     :: !String
  } deriving (Eq, Show)

-- The mark is FUSED onto the stage, so the two unreachable states
--   (FutRegistered, Just p)  and  (FutExpired, Nothing)  are unspellable.
-- FutRegistered carries no mark; FutExpired always carries one; FutActive may or may not.
data Settlement = Settlement { settlePrice :: !Price, settleDate :: !Day }
  deriving (Eq, Show)

data FutStage
  = FutRegistered                -- recorded, never traded: no holders, no mark
  | FutActive (Maybe Settlement) -- trading; Nothing = not yet settled, Just s = marked
  | FutExpired Settlement        -- final mark ALWAYS present; absorbing
  deriving (Eq, Show)

newtype FutStatus = FutStatus { fusStage :: FutStage } deriving (Eq, Show)

-- last_settlement_price/_date are PROJECTIONS of the mark the stage carries,
-- total over all three stages -- not independent fields.
settlement :: FutStage -> Maybe Settlement
settlement FutRegistered  = Nothing
settlement (FutActive m)  = m
settlement (FutExpired s) = Just s

settlementPrice :: FutStatus -> Maybe Price
settlementPrice = fmap settlePrice . settlement . fusStage
settlementDate  :: FutStatus -> Maybe Day
settlementDate  = fmap settleDate  . settlement . fusStage

-- The coarse rank REGISTERED < ACTIVE < EXPIRED: the one fact futApplyDelta needs
-- to reject a regressing stage. It never regresses.
stageRank :: FutStage -> Int
stageRank FutRegistered  = 0
stageRank (FutActive _)  = 1
stageRank (FutExpired _) = 2

-- Rank alone cannot separate FutExpired 105 from FutExpired 110, so a strict new<cur
-- guard would re-admit a second Expire. isExpired is the precise absorbing test.
isExpired :: FutStage -> Bool
isExpired (FutExpired _) = True
isExpired _              = False

futRegisteredStatus :: FutStatus
futRegisteredStatus = FutStatus FutRegistered
\end{lstlisting}

{\raggedright\noindent\emph{Evaluations.}
\code{settlementPrice (FutStatus FutRegistered) = Nothing};
\code{settlementPrice (FutStatus (FutExpired (Settlement (Price 105) (Day 3)))) =
Just (Price 105)}. The states \code{(FutRegistered, Just p)} and
\code{(FutExpired, Nothing)} cannot be written.\par}

%==============================================================================
\subsection{First trade}
\label{sec:trade}
%==============================================================================

A buys 10 contracts from B at 100. The trade emits one \code{FutStateDelta}:

\begin{itemize}
\item \textbf{Shared.} \code{REGISTERED} $\to$ \code{ACTIVE} with no mark (\code{FutActive Nothing}); the
first trade activates the unit. Any existing mark is preserved --- a trade never wipes the shared price.
\item \textbf{Per-position.} The contract quantity moves from seller to buyer:
$\Delta\netq(A)=+10$, $\Delta\netq(B)=-10$. Each leg sets $\ac\mathrel{+}=-\Delta\cdot p\cdot\mult$:
\[
\Delta\ac(A) = -(+10)\cdot 100\cdot 50 = -50000,\qquad
\Delta\ac(B) = -(-10)\cdot 100\cdot 50 = +50000.
\]
\item \textbf{Cash.} None: entering a futures position pays no notional.
\end{itemize}

\noindent Conservation: $\sum\Delta\netq = 0$, $\sum\Delta\ac = 0$, $\sum\mathrm{VM}=0$. State after
T1: $\netq=(10,-10,0)$, $\ac=(-50000,+50000,0)$. A and B cross from $\mathrm{None}$ to $\mathrm{Some}$
on first touch; C is still $\mathrm{None}$.

\paragraph{Positive quantity.} A trade leg carries a strictly positive quantity $q>0$, parsed at the
boundary (\code{mkPosQty}). $q\le 0$ is rejected (\code{NonPositiveQty}) and never recorded. $q=0$
would promote \code{REGISTERED}$\to$\code{ACTIVE} and create $\mathrm{Some}$-flat rows for two
never-held wallets, collapsing the never-held/held-flat distinction (\S\ref{sec:close}). $q<0$ would
swap buyer and seller, violating the move primitive's positivity. The sign of the position change is
supplied by the leg --- buyer $+q$, seller $-q$ --- not by $q$.

The per-event delta, its conserved image, and the conservation witness follow. The mental model:
totalling a delta is addition in a monoid, ``conserving'' means the total is zero, and the witness
type makes an unconserved delta impossible to apply.

\begin{lstlisting}[language=Haskell]
-- FutPos per (wallet, unit). ac is per (w,u): two wallets holding the same
-- contract carry different ac -- the field that makes intraday margin correct.
data FutPos = FutPos { fpNetQty :: !Qty, fpAccumCost :: !Cash }
  deriving (Eq, Show)

futZeroP :: FutPos               -- flat row created on first touch; held-and-flat baseline
futZeroP = FutPos mempty mempty

-- Conservation is a monoid identity. Totalling a delta is a monoid homomorphism
-- -- a map that respects (<>) and mempty -- into Conserved; "conserving" means the
-- image is mempty. The zero-holder case is the empty fold, mempty (C9); we SUM
-- deltas, never apportion, so no divide-by-holder-count bug can arise.
data Conserved = Conserved !Qty !Cash !Cash deriving (Eq, Show)  -- (sum net, sum ac, sum cash)
instance Semigroup Conserved where
  Conserved a b c <> Conserved a' b' c' = Conserved (a<>a') (b<>b') (c<>c')
instance Monoid Conserved where mempty = Conserved mempty mempty mempty

data FutRowDelta = FutRowDelta { rdNetQty :: !Qty, rdAc :: !Cash } deriving (Eq, Show)

-- One event's proposed change: at most one shared FutStage write, per-holder additive
-- row deltas, per-holder cash legs. Applied in full or not at all (C3).
data FutStateDelta = FutStateDelta
  { fsdUnit :: UnitId, fsdStage :: Maybe FutStage
  , fsdRows :: Map WalletId FutRowDelta, fsdCash :: Map WalletId Cash } deriving (Show)

futNetDelta :: FutStateDelta -> Conserved
futNetDelta sd =
       foldMap (\(FutRowDelta nq ac) -> Conserved nq ac mempty) (fsdRows sd)
    <> foldMap (\c                   -> Conserved mempty mempty c) (fsdCash sd)

-- The C2 witness, ABSTRACT: futValidate is the only constructor, so an unconserved
-- delta cannot reach futApplyDelta. This is the futures engine's OWN conservation
-- discharge, distinct from the core (which has no validate gate); the validated
-- delta projects onto one conserving core Transaction. Conservation is a value-level
-- fact (a sum of exact integers), so this is a parse boundary, not a free type fact.
newtype FutValidDelta = FutValidDelta FutStateDelta

futValidate :: FutStateDelta -> Either FutLedgerError FutValidDelta
futValidate sd | net == mempty = Right (FutValidDelta sd)
               | otherwise     = Left (FutNotConserved (fsdUnit sd) net)
  where net = futNetDelta sd

-- A trade activates without wiping any existing mark, and sets each leg's
-- ac += -(delta signed_qty) * p * m (buyer +q, seller -q). No cash: entry pays no notional.
activateTrade :: FutStage -> FutStage
activateTrade cur = FutActive (settlement cur)

tradeDelta :: UnitId -> FutStage -> Integer
           -> WalletId -> WalletId -> PosQty -> Price -> FutStateDelta
tradeDelta u cur m buyer seller pq p = FutStateDelta
  { fsdUnit = u, fsdStage = Just (activateTrade cur)
  , fsdRows = Map.fromList
      [ (buyer , FutRowDelta q        (cashNeg (futMark q p m)))
      , (seller, FutRowDelta (qneg q) (futMark q p m)) ]
  , fsdCash = Map.empty }
  where q = unPosQty pq
\end{lstlisting}

%==============================================================================
\subsection{Daily variation-margin settlement: the mechanism and the clean case (day 1)}
\label{sec:settle-mech}
%==============================================================================

Settlement pays each holder the day's mark-to-market in cash and re-bases the position at the new
price. Formally: settlement at price $S$, multiplier $\mult$, over the current holders. For each
holder the reset target is $\mathrm{target}=-\netq\cdot S\cdot\mult$, the row delta is
$\Delta\ac=\mathrm{target}-\ac$, and the cash leg is
\[
\boxed{\;\mathrm{VM}(w) \;=\; -\,\Delta\ac(w) \;=\; \netq(w)\cdot S\cdot \mult + \ac(w)\;}
\]
The cash leg is the mirror of the change in accumulated cost. Therefore
$\sum_w\mathrm{VM} = -\sum_w\Delta\ac$, and variation-margin zero-sum is the \emph{same} fact as \ac{}
conservation, not a separate reconciliation. The reset discards the entry prices: with
$\ac\leftarrow-\netq\cdot S\cdot\mult$, the next day's margin needs only $S$ and the position. Each
daily settlement economically closes and reopens the position at the new price (the resettable-forward
property).

Settlement is one atomic \code{FutStateDelta} touching both layers:
\begin{itemize}
\item \textbf{Shared}, one write: \emph{UnitStatus}$[\unit]$ acquires the mark $(S,d)$. The
\emph{coarse rank} $\code{REGISTERED}<\code{ACTIVE}<\code{EXPIRED}$ is unchanged: the unit stays
\code{ACTIVE}. The \emph{embedded mark} updates: the stage becomes
\code{FutActive (Just (Settlement $S$ $d$))}. The settlement price enters the log only as the
\code{FSettleVM} event. \code{futReplay} rebuilds the mark, and the fold stays pure at the
boundary: the shared value is a projection of the log, never an off-log write.
\item \textbf{Per-position}, a fan-out over the current holders: each row's \ac{} is reset to
$\mathrm{target}$, and each holder gets the cash leg $\mathrm{VM}(w)$.
\end{itemize}

\noindent The cash legs are emitted over the holders; the clearinghouse leg is the residual that
balances them. Since the holder legs already sum to zero, CH's leg is zero throughout this example and
no CH cash row materialises. This is the special case of a \emph{closed} holder set $\{A,B,C\}$, where
$\sum_w\netq(w)=0$ at every step. In the Ledger's intended use --- an internal firm system of record
facing an external clearinghouse --- the firm is generally \emph{not} flat. When $\sum_w\netq(w)\ne 0$
over the firm's holders, the CH leg is the firm's nonzero net variation margin: the boundary figure
reconciled against the clearing broker's daily statement.

\paragraph{Day 1: $S_1=102$.} Holders A and B; C has no row and is untouched --- the contract settles
vacuously over C; the shared price still updates (C9). With pre-settle $\ac=(-50000,+50000)$ and
$\netq=(10,-10)$:
\[
\mathrm{VM}(A) = 10\cdot 102\cdot 50 - 50000 = +1000,\qquad
\mathrm{VM}(B) = -10\cdot 102\cdot 50 + 50000 = -1000.
\]
Reset: $\ac(A)\leftarrow -51000$, $\ac(B)\leftarrow +51000$. Conservation: $\sum\Delta\netq=0$,
$\sum\Delta\ac = 0$, $\sum\mathrm{VM}=0$. State: $\netq=(10,-10,0)$, $\ac=(-51000,+51000,0)$.

The fan-out, shared by settlement and expiry, is one pure function of the current holders. Per
holder it does two things: reset \ac{} to the new-price baseline, and pay the mirror of that reset
as cash.

\begin{lstlisting}[language=Haskell]
-- Shared by FSettleVM and FExpire. For each holder at price S, multiplier m:
--   delta ac = (-(net*S*m)) - ac      VM = -(delta ac) = net*S*m + ac
-- so VM = -(delta ac) exactly: cash conservation IS ac conservation, not a separate
-- reconciliation. A flat holder gives delta ac = 0, VM = 0; no holders -> empty fan-out.
settlementFanout :: Price -> Integer -> [(WalletId, FutPos)]
                 -> (Map WalletId FutRowDelta, Map WalletId Cash)
settlementFanout s m holders =
    ( Map.fromList [ (w, FutRowDelta mempty (deltaAc ps)) | (w, ps) <- holders ]
    , Map.fromList [ (w, cashNeg (deltaAc ps))            | (w, ps) <- holders ] )
  where deltaAc ps = cashNeg (futMark (fpNetQty ps) s m) <> cashNeg (fpAccumCost ps)
\end{lstlisting}

%==============================================================================
\subsection{A subsequent, intraday trade (day 2)}
%==============================================================================

C buys 4 contracts from A at 103, after the day-1 settle and before the day-2 settle --- A trades
intraday. One \code{FutStateDelta}, stage already \code{FutActive} (mark preserved):
\[
\Delta\netq(C)=+4,\quad \Delta\netq(A)=-4;\qquad
\Delta\ac(C) = -20600,\quad \Delta\ac(A) = +20600.
\]
Conservation: $\sum\Delta\netq=0$, $\sum\Delta\ac=0$, $\sum\mathrm{VM}=0$. State after T2:
$\netq=(6,-10,4)$, $\ac=(-30400,+51000,-20600)$. C crosses from $\mathrm{None}$ to $\mathrm{Some}$ on
first touch.

%==============================================================================
\subsection{Settlement under intraday trading (day 2): the anchor}
\label{sec:anchor}
%==============================================================================

{\raggedright
The day-2 settle at $S_2=101$ runs the same fan-out over holders A, B, C:
\[
\begin{aligned}
\mathrm{VM}(A) &= 6\cdot 101\cdot 50 - 30400 = -100,\\
\mathrm{VM}(B) &= -10\cdot 101\cdot 50 + 51000 = +500,\\
\mathrm{VM}(C) &= 4\cdot 101\cdot 50 - 20600 = -400.
\end{aligned}
\]
Conservation: $\sum\Delta\netq=0$, $\sum\mathrm{VM} = -100+500-400 = 0$. Reset
$\ac\leftarrow(-30300,+50500,-20200)$ with $\sum\Delta\ac=0$. State: $\netq=(6,-10,4)$,
$\ac=(-30300,+50500,-20200)$.\par}

\begin{principle}[The intraday-margin subtlety forces stored per-position \ac{} (C11)]
A's day-2 variation margin is $-100$, not the naive $\netq\cdot(S_2-S_1)\cdot\mult =
6\cdot(101-102)\cdot 50 = -300$. A sold 4 contracts at 103 intraday, one point above the prior mark
102, gaining $4\cdot(+1)\cdot 50=+200$, which offsets the $-300$ mark loss on the held position. Only
the stored per-wallet \ac{}, which absorbed the intraday trade, yields the correct $-100$. A
price-derived model gives the wrong cash. Accumulated cost is therefore per-position stored state with
a single canonical writer (C11), not a consequence read off the shared price.
\end{principle}

\noindent The three settlement facts, stated once:
\begin{enumerate}
\item \emph{Settlement is a state update split by layer.} The shared part ---
\code{last\_settlement\_price}, \code{last\_settlement\_date} --- is one write on
\emph{UnitStatus}$[\unit]$, one value per contract. The per-position part --- the \ac{} reset and the
cash leg --- is a fan-out over the current holders on \emph{PositionState}$[w,\unit]$. One atomic
\code{FutStateDelta} touches both.
\item \emph{One atomic event that fans out, not a derived consequence of the price.} The cash leg is
real daily money moving between longs and shorts through CH. Every cash move is a recorded,
conservation-bearing event, so a per-holder pass is unavoidable. The \ac{} reset rides in the same
atomic delta. The fan-out is forced by the cash, not chosen for the bookkeeping.
\item \emph{Price only in shared state, consequence only in per-wallet state.} The consequence of
the shared price --- the \ac{} reset and the variation-margin cash --- lives only in
\emph{PositionState}$[w,\unit]$ and the move stream.
\end{enumerate}

\begin{remark}[Direction reversal needs no special case]
The accumulated-cost rule handles a long-to-short reversal in one trade with no branch. From
$\netq=+5$, $\ac=-500$ (entered at 100, $\mult=1$), a single sell of 8 at 105 sets
$\Delta\ac=-(-8)\cdot 105\cdot 1=+840$, so $\ac=-500+840=+340$ and $\netq=-3$. The economic value at
105 is $340+(-3)\cdot 105 = 25$, exactly the gain from closing the long 5 at \$5 per contract; the new
short 3 marks forward from here with no extra bookkeeping. The same arithmetic that accumulates a
position unwinds and reverses it.
\end{remark}

%==============================================================================
\subsection{Closing a position to flat}
\label{sec:close}
%==============================================================================

B buys 4 contracts from C at 101; C goes flat. One \code{FutStateDelta}:
\[
\Delta\netq(B)=+4,\ \Delta\netq(C)=-4;\qquad \Delta\ac(B)=-20200,\ \Delta\ac(C)=+20200.
\]
Conservation: $\sum\Delta\netq=0$, $\sum\Delta\ac=0$, $\sum\mathrm{VM}=0$. State: $\netq=(6,-6,0)$,
$\ac=(-30300,+30300,0)$. C is now flat. Its row is \emph{retained}, $\mathrm{Some}$ at $(0,0)$ --- not
deleted, not $\mathrm{None}$. The two readings are distinct and both load-bearing. $\mathrm{None}$
means never held. $\mathrm{Some}$-flat means held and now flat; it is read after close-out for tax
reporting, wash-sale lookback, and trade reconstruction (C1(a)). The carrier is monotone: once a row
is created it is never removed (C1(b)). No row deleter exists, because the ledger is abstract.

%==============================================================================
\subsection{Expiry, final settlement, and the Close}
%==============================================================================

Expiry is the settlement fan-out of \S\ref{sec:settle-mech} composed with a stage write to
\code{EXPIRED}. Like a settle, it acts only on an \code{ACTIVE} unit, never a never-traded one
(\S\ref{sec:invariants}). The final mark $S=105$ fans out over the remaining holders A and B (C is
flat; its row is touched to no effect):
\[
\mathrm{VM}(A) = 6\cdot 105\cdot 50 - 30300 = +1200,\qquad
\mathrm{VM}(B) = -6\cdot 105\cdot 50 + 30300 = -1200.
\]
Conservation: $\sum\Delta\netq=0$, $\sum\mathrm{VM}=0$; $\ac\leftarrow(-31500,+31500,0)$ with
$\sum\Delta\ac=0$. The coarse rank advances to \code{EXPIRED}: the stage becomes
\code{FutExpired (Settlement 105 (Day 3))}; the final mark is always present.

\paragraph{Cash settlement and the Close.} No unit is delivered. The economic result is the
cumulative variation margin already moved. The \emph{Close} step extinguishes every holder's position
against the clearinghouse, flattening $\netq$ and \ac{} by an additive negation of each row, with no
further cash. Wallet A held $+6$ and B held $-6$; closing them returns the contracts to CH:
\[
\Delta\netq(A)=-6,\ \Delta\netq(B)=+6;\qquad \Delta\ac(A)=+31500,\ \Delta\ac(B)=-31500.
\]
The clearinghouse legs are the residual of the holder legs --- zero, since the holder legs balance ---
so no CH row materialises. Conservation: $\sum\Delta\netq=0$, $\sum\Delta\ac=0$, $\sum\mathrm{VM}=0$.
State after Close: $\netq=(0,0,0)$, $\ac=(0,0,0)$; rows retained at zero (monotone carrier). Close
carries no stage write, so it is the one event still admissible on an \code{EXPIRED} unit.

\paragraph{Physical settlement.} The variant differs only at the Close. Instead of flattening against
CH for cash, the contract converts to delivery. The deliverable unit and the cash payment move
together as one atomic transaction (delivery versus payment; the settlement interface is
\S\ref{sec:settlement}). The future's $\netq$ goes to zero, and the holder's wallet gains the
deliverable. The cash leg is struck at the \emph{final settlement price} $S$, not the trade price.
Expiry has already marked every position to $S$ through variation margin, so delivery exchanges the
deliverable for $\netq\cdot S\cdot\mult$ and moves no further mark-to-market. Pricing the delivery at
any other figure would double-count the variation margin already settled. Per-unit conservation holds
on each unit separately --- the future, the deliverable, the cash --- and the cross-unit delivery is
one all-or-nothing transaction. The worked example is cash-settled; the figures above are unchanged by
this variant.

\begin{theorem}[Closing identity: cash ties to economic PnL]
Cumulative variation margin per wallet equals the wallet's economic profit and loss. Here
$A=+1000-100+1200=+2100$, $B=-1000+500-1200=-1700$, $C=-400$, $\text{CH}=0$, summing to $0$; and
$A=4\cdot(103-100)\cdot 50+6\cdot(105-100)\cdot 50=+2100$,
$B=4\cdot(100-101)\cdot 50+6\cdot(100-105)\cdot 50=-1700$, $C=4\cdot(101-103)\cdot 50=-400$.
\end{theorem}

The terminal flatten delta has no stage write and no cash leg. It conserves for free: each conserved
field already sums to zero over the holders, so the row-by-row negation sums to zero too.

\begin{lstlisting}[language=Haskell]
-- Terminal flatten: NO stage write (stays EXPIRED, so it does not breach the
-- absorbing rule), NO cash; every row driven to zero by additive negation. It
-- conserves because net and ac each already sum to zero, so their negations do too.
closeDelta :: UnitId -> [(WalletId, FutPos)] -> FutStateDelta
closeDelta u holders = FutStateDelta
  { fsdUnit = u, fsdStage = Nothing
  , fsdRows = Map.fromList
      [ (w, FutRowDelta (qneg (fpNetQty ps)) (cashNeg (fpAccumCost ps)))
      | (w, ps) <- holders ]
  , fsdCash = Map.empty }
\end{lstlisting}

%==============================================================================
\subsection{After expiry: state versus derivation}
%==============================================================================

After the Close, every row is retained at zero: \emph{PositionState}$[w,\unit]=\mathrm{Some}$ flat for
A, B, C. \emph{UnitStatus}$[\unit]$ holds stage \code{EXPIRED} and final mark $105$.
\emph{ProductTerms}$[\unit]$ is unchanged --- immutable.

\code{first\_touch\_date} is \emph{not} stored: it is derived from the event log on demand. Caching it
in \emph{PositionState} would make a replay disagree with itself under a back-dated correction: the
cached value would reflect the pre-correction order while the fold reflects the corrected one. The
rule is general. What the fold over the log determines is derived, not stored. Only what the fold
cannot reconstruct from prior events is state.

%==============================================================================
\subsection{Invariants threaded through the life}
\label{sec:futures-invariants}
%==============================================================================

\begin{invariant}[Per-event conservation, C2]
Every \code{FutStateDelta} satisfies $\sum_w\Delta\netq=0$, $\sum_w\Delta\ac=0$, $\sum_w\mathrm{VM}=0$,
discharged by \code{futValidate} before application. Induction over the stream lifts the per-event
identity to the global invariant (P1). The vacuous case --- a unit with no holders --- is the empty
sum, zero (C9); conservation sums deltas and never divides by a holder count.
\end{invariant}

\begin{invariant}[Deterministic replay as a fold, C1(b), P8]
\code{futReplay} is a Kleisli fold (an error-stopping fold whose combining step halts at the first failure) of the event stream over the ledger:
$\code{futReplay}(xs\mathbin{<\!\!>}ys)=\code{futReplay}\,xs \mathbin{>\!=\!>} \code{futReplay}\,ys$ ---
replaying a concatenation equals composing the replays. Checkpoint independence is a consequence of
this law, not a test; the monotone carrier keeps the key set stable across cuts. The stream is assumed
\emph{unique} --- de-duplicated at the ingestion boundary by an event key --- because the two event
kinds differ in idempotency: a duplicated \code{FSettleVM} at a fixed mark is inert ($\Delta\ac=0$,
$\mathrm{VM}=0$), whereas a duplicated \code{FTrade} accumulates a second position delta. Duplicate
suppression is a boundary obligation, not a ledger function.
\end{invariant}

\begin{invariant}[Monotone, absorbing stage]
The coarse rank $\code{REGISTERED}<\code{ACTIVE}<\code{EXPIRED}$ never regresses; a proposed stage of
strictly lower rank is rejected (\code{StageRegression}). Settle and expire act only on an
\code{ACTIVE} unit; both reject a \code{REGISTERED} (never-traded) unit (\code{NotActive}), which has
no holders and no mark slot --- so a never-traded unit is never silently promoted, and \code{EXPIRED}
is reachable only from \code{ACTIVE}: the chain is
\code{REGISTERED}$\to$\code{ACTIVE}$\to$\code{EXPIRED} with no skips. \code{EXPIRED} is
\emph{absorbing}: the rank guard alone is too weak --- a second expiry has equal rank, and $2<2$ is
false --- so an explicit absorbing test rejects every stage-writing delta on an expired unit;
\code{FClose}, which carries no stage write, is the one event still admissible. Re-settling at the same
price before expiry is a no-op: with $\ac=-\netq\cdot S\cdot\mult$ already, $\Delta\ac=0$ and
$\mathrm{VM}=0$.
\end{invariant}

The event dispatcher, the abstract ledger, and the replay fold realise these invariants. One event
flows through three total functions: \code{futHandle} decides what should happen, \code{futValidate}
certifies conservation, and \code{futApplyDelta} writes. The stage and positivity guards live once in
\code{futHandle}. The monotone carrier is the absence of any deleter; the absorbing rule is the
\code{isExpired} test in \code{futApplyDelta}.

\begin{lstlisting}[language=Haskell]
-- The four event classes. FTrade carries a RAW Qty (the unparsed input of the free
-- monoid of events); futHandle parses it once into PosQty. FClose is the one event
-- with no stage write, so it alone is admissible on an EXPIRED unit.
data FutEvent
  = FTrade    UnitId WalletId WalletId Qty Price   -- buyer, seller, qty, price
  | FSettleVM UnitId Price Day                      -- settlement price, date
  | FExpire   UnitId Price Day                      -- final price, date
  | FClose    UnitId                                -- terminal flatten to zero
  deriving (Eq, Show)

futEventUnit :: FutEvent -> UnitId
futEventUnit (FTrade u _ _ _ _) = u
futEventUnit (FSettleVM u _ _)  = u
futEventUnit (FExpire u _ _)    = u
futEventUnit (FClose u)         = u

-- The dispatcher: a pure function of (event, current ledger), hence deterministic.
-- The stage/positivity guards live here once: self-trade (G3); EXPIRED absorbing
-- (G2); positivity parsed (G5); settle/expire reject REGISTERED, never silently
-- promoting a never-traded unit (G4); close only after expiry (G1).
futHandle :: FutEvent -> FutLedger -> Either FutLedgerError FutStateDelta
futHandle ev l = case Map.lookup u (flUnits l) of
  Nothing          -> Left (FutUnknownUnit u)
  Just (terms, us) ->
    let m = ptMultiplier terms; cur = fusStage us in case ev of
      FTrade _ buyer seller q p
        | buyer == seller -> Left (SelfTrade u buyer)            -- G3
        | isExpired cur   -> Left (UnitExpired u)                -- G2
        | otherwise       -> case mkPosQty q of                  -- G5: parse once
            Nothing -> Left (NonPositiveQty u q)
            Just pq -> Right (tradeDelta u cur m buyer seller pq p)
      FSettleVM _ s d -> case cur of
        FutExpired _  -> Left (UnitExpired u)                     -- G2
        FutRegistered -> Left (NotActive u)                       -- G4
        FutActive _   -> let (rows, cash) = settlementFanout s m (futHoldersOf u l)
                         in Right (FutStateDelta u (Just (FutActive (Just (Settlement s d)))) rows cash)
      FExpire _ s d -> case cur of
        FutExpired _  -> Left (UnitExpired u)                     -- G2
        FutRegistered -> Left (NotActive u)                       -- G4 (symmetric)
        FutActive _   -> let (rows, cash) = settlementFanout s m (futHoldersOf u l)
                         in Right (FutStateDelta u (Just (FutExpired (Settlement s d))) rows cash)
      FClose _ -> case cur of
        FutExpired _ -> Right (closeDelta u (futHoldersOf u l))   -- G1
        _            -> Left (NotExpired u)
  where u = futEventUnit ev

-- The ledger, ABSTRACT: no field setter or row/status/cash deleter is exported, so
-- futApplyDelta (insert/update only) is the sole mutator and that absence IS the
-- monotone-carrier discipline (C1). Terms and status are FUSED into one map, so the
-- desynchronised state (one map has u, the other does not) is unrepresentable.
data FutLedger = FutLedger
  { flUnits :: Map UnitId (FutTerms, FutStatus)
  , flPos   :: Map (WalletId, UnitId) FutPos
  , flCash  :: Map WalletId Cash }

data FutLedgerError
  = FutReRegistration UnitId | FutUnknownUnit UnitId
  | StageRegression UnitId FutStage FutStage | FutNotConserved UnitId Conserved
  | UnitExpired UnitId | NotActive UnitId | NotExpired UnitId
  | SelfTrade UnitId WalletId | NonPositiveQty UnitId Qty
  deriving (Eq, Show)

futEmptyLedger :: FutLedger
futEmptyLedger = FutLedger Map.empty Map.empty Map.empty

-- Co-registers terms AND status as one tuple. Re-registration is a hard error (C10).
futRegister :: UnitId -> FutTerms -> FutLedger -> Either FutLedgerError FutLedger
futRegister u terms l
  | u `Map.member` flUnits l = Left (FutReRegistration u)
  | otherwise = Right l { flUnits = Map.insert u (terms, futRegisteredStatus) (flUnits l) }

-- Apply a validated delta: total, atomic, monotone. EXPIRED rejects ANY stage write
-- (absorbing); a bare new<=cur would wrongly reject the ACTIVE->ACTIVE daily
-- re-settle, so the precise rule is isExpired plus a rank-regression guard. FClose's
-- fsdStage = Nothing bypasses both. Rows and cash are inserted/updated, never removed.
futApplyDelta :: FutValidDelta -> FutLedger -> Either FutLedgerError FutLedger
futApplyDelta (FutValidDelta sd) l
  | not (u `Map.member` flUnits l) = Left (FutUnknownUnit u)
  | otherwise = case fsdStage sd of
      Nothing  -> Right written
      Just new
        | isExpired cur                 -> Left (UnitExpired u)
        | stageRank new < stageRank cur -> Left (StageRegression u cur new)
        | otherwise -> Right written
            { flUnits = Map.adjust (\(t,_) -> (t, FutStatus new)) u (flUnits l) }
  where
    u   = fsdUnit sd
    cur = maybe FutRegistered (fusStage . snd) (Map.lookup u (flUnits l))
    written = l
      { flPos  = Map.foldrWithKey applyRow (flPos l) (fsdRows sd)
      , flCash = Map.foldrWithKey (\w c -> Map.insertWith (<>) w c)
                                  (flCash l) (fsdCash sd) }
    applyRow w (FutRowDelta dnq dac) mp =
      let key = (w, u); p0 = Map.findWithDefault futZeroP key mp   -- first touch -> flat
      in Map.insert key p0 { fpNetQty    = fpNetQty    p0 <> dnq
                           , fpAccumCost = fpAccumCost p0 <> dac } mp

-- One event end to end: read state -> discharge conservation -> apply. Each futStep
-- projects onto one core Transaction recorded through the core applyTx; futStep is
-- the futures engine's narration of that projection, not a second core door.
futStep :: FutEvent -> FutLedger -> Either FutLedgerError FutLedger
futStep ev l = futHandle ev l >>= futValidate >>= \vd -> futApplyDelta vd l

-- futReplay is a Kleisli fold -- a fold whose step composes in the Either monad. Its
-- homomorphism law  futReplay (xs <> ys) = futReplay xs >=> futReplay ys  makes
-- checkpoint independence a CONSEQUENCE, not a test. The stream is assumed
-- de-duplicated at ingestion: a duplicate FSettleVM at a fixed mark is inert, but a
-- duplicate FTrade accumulates a second delta.
futReplay :: [FutEvent] -> FutLedger -> Either FutLedgerError FutLedger
futReplay evs l0 = foldM (flip futStep) l0 evs

-- Total accessors. futPosition is the C1 Option accessor: Nothing = never held,
-- Just futZeroP = held and now flat -- never collapsed.
futProductTerms :: FutLedger -> UnitId -> Maybe FutTerms
futProductTerms l u = fst <$> Map.lookup u (flUnits l)
futUnitStatus :: FutLedger -> UnitId -> Maybe FutStatus
futUnitStatus l u = snd <$> Map.lookup u (flUnits l)
futPosition :: FutLedger -> WalletId -> UnitId -> Maybe FutPos
futPosition l w u = Map.lookup (w, u) (flPos l)
futCashOf :: FutLedger -> WalletId -> Cash
futCashOf l w = Map.findWithDefault mempty w (flCash l)
futHoldersOf :: UnitId -> FutLedger -> [(WalletId, FutPos)]
futHoldersOf u l = [ (w, ps) | ((w, u'), ps) <- Map.toList (flPos l), u' == u ]
\end{lstlisting}

%==============================================================================
\subsection{Escalations}
%==============================================================================

The two escalations that follow are carried to the register in \S\ref{sec:conclusion}.

\paragraph{FE1 --- Fan-out cost at scale.} A settlement writes one \emph{PositionState} row and emits
one cash leg per open holder, every day. Across $\sim\!10^{6}$ contracts each with up to thousands of
holders, daily settlement is $O(\text{open positions})$ writes and cash legs --- a large daily
fan-out and write amplification. This is intrinsic to per-wallet variation margin. Mitigations
(batching, snapshotting) trade against the conservation-per-event and single-writer disciplines and
must be shown to preserve C2 and C11. The cost stands. \emph{Owner:} performance and risk engineering.

\paragraph{FE2 --- The derived-consequence alternative, declined.} Storing only the shared
\code{last\_\allowbreak settlement\_\allowbreak price} and computing each wallet's mark lazily saves
nothing on the dominant cost: the cash leg is per-wallet and unavoidable (FE1). It cannot produce the
correct variation margin for an intraday trader without reconstructing \ac{} anyway
(\S\ref{sec:anchor}). It would also split the single-writer discipline for \ac{} (C11).

\subsection*{Key Invariants and Consequences}

Valuation and lifecycle processing are projections of state; the following hold once
conservation is granted. Each is catalogued in \S\ref{sec:invariants}.

\begin{itemize}[nosep]
\item \textbf{State-sufficiency.} Portfolio value depends only on current state---wallet
balances, unit state (including embedded lifecycle state), and the current price
vector---not on transaction history. Valuation is the linear functional
$V_t=\sum_u w_t(u)\,P_t(u)$. (\S\ref{sec:valuation}.)
\item \textbf{PnL path-independence (P10).} $\mathrm{PnL}=V_{t_1}-V_{t_0}$, independent of
the intermediate transactions. Dual valuation applies two price vectors to one state.
(\S\ref{sec:valuation}, Theorem~\ref{inv:P10}.)
\item \textbf{Lifecycle value invariance.} Across deterministic lifecycle
events---coupons, dividends, resets, redemptions---total value under a given price vector
is invariant: the instrument's price jump is offset by the explicit cash move. Under
optionality the invariance holds for intrinsic value only. (\S\ref{sec:intro}; \S\ref{sec:lifecycle}.)
\item \textbf{Contracts as move generators.} A smart contract is a deterministic program
emitting sequences of moves; because every lifecycle effect is a move, conservation P1 is
preserved with no separate validation. (\S\ref{sec:contracts}.)
\item \textbf{Lifecycle idempotency (P6).} Applying the same lifecycle event to the same
unit state produces no additional moves. (\S\ref{sec:lifecycle}, Invariant~\ref{inv:P6}.)
\item \textbf{Lifecycle purity (P9).} For fixed inputs, lifecycle functions produce
identical outputs and no effect beyond the return, so valuation and time travel are
deterministic. (\S\ref{sec:lifecycle}.)
\end{itemize}


